\documentclass[12pt,a4paper]{ctexart}
\usepackage[top=2.5cm,bottom=2.5cm,left=2.5cm,right=2.5cm]{geometry}
\usepackage{amsmath,amssymb,amsthm,physics}
\usepackage{graphicx,float,booktabs,array,caption,multirow}
\usepackage{hyperref,xcolor,enumitem,mathtools}
\usepackage[numbers,sort&compress]{natbib}
\hypersetup{colorlinks=true,linkcolor=blue,citecolor=blue,urlcolor=blue}
\theoremstyle{definition}
\newtheorem{definition}{定义}[section]
\newtheorem{theorem}{定理}[section]
\newtheorem{remark}{注记}[section]

\title{\heiti 格点QCD中的衰变道：\\
从衰变常数到L\"uscher有限体积方法的完整解析}
\author{基于本库全部相关文献与教材的综合解析}
\date{\today}

\begin{document}
\maketitle

\begin{abstract}
\noindent
强子的衰变——从最轻的$\pi$介子和$K$介子的轻子衰变，到$\rho$介子到$\pi\pi$的强衰变，再到重味$B$和$D$介子的半轻子和非轻子衰变——是格点QCD中最具挑战性也最富物理成果的研究领域之一。与强子质量（可通过两点关联函数的指数衰减简单提取）不同，衰变过程涉及Minkowski时空中的实时动力学——未态粒子之间的强相互作用使得欧几里得格点计算面临根本性的理论困难。尽管如此，过去三十年已有三类互补的理论工具在格点上被发展和完善：\textbf{衰变常数}（通过真空与单强子态的矩阵元直接提取）、\textbf{L\"uscher有限体积方法}（将有限体积中的双粒子能谱映射到无穷大体积中的散射相移和共振参数）、以及\textbf{有效弱Hamiltonian方法}（通过算符乘积展开分离短程和长程物理）。

本文基于本库中的核心教材（Gattringer与Lang的\textit{Quantum Chromodynamics on the Lattice}——第11章完整推导了衰变常数、$\pi\pi$散射的L\"uscher方法、$K\to\pi\pi$和弱衰变的理论框架；Gupta的\textit{Introduction to Lattice QCD}——第14章和第19章提供了散射和衰变的实践指南；Peskin与Schroeder的\textit{An Introduction to Quantum Field Theory}——Part III的QCD背景）以及约15篇研究论文，从以下六个层次系统解析格点QCD中的衰变道：

\textbf{第一层——衰变常数的第一性原理计算（§2）：}
从PCAC（部分守恒轴矢流）关系出发，推导$\pi$介子衰变常数$F_\pi$和$K$介子衰变常数$F_K$的格点提取公式：$F_\pi$从轴矢流关联函数的大$t$行为$\langle A_4 A_4\rangle \sim (M_\pi F_\pi^2/2) e^{-M_\pi t}$中提取。详述粲介子衰变常数$f_{D_{(s)}}$和$f_{D_{(s)}^*}$的现代格点计算——包括Tadpole改进Clover系综上$f_{D_s} = 249.9(0.5)$~MeV的精确结果~\cite{CharmedClover}，以及这些衰变常数在确定CKM矩阵元$|V_{cd}|$和$|V_{cs}|$中的关键角色。

\textbf{第二层——强子衰变与L\"uscher有限体积方法（§3）：}
这是连接有限体积欧几里得格点与无穷大体积Minkowski散射/衰变物理的核心理论工具。系统阐述L\"uscher公式的推导——从一维周期边界条件下平面波函数的相移匹配条件$e^{ikL+2i\delta(k)} = 1$出发，到四维QCD中能量本征值$W_n(L)$与散射相移$\delta(k)$之间的泛函关系：$\delta(k) = \phi(q) \mod \pi$，$\tan(-\phi(q)) = q\pi^{3/2}/Z_{00}(1;q^2)$。详细讨论双粒子基态能量与$s$波散射长度$a_0$的关系（$W = 2M_\pi - (4\pi a_0/M_\pi L^3)[1 + c_1 a_0/L + c_2 a_0^2/L^2]$），以及$\rho\to\pi\pi$衰变作为典范非静止框架的推广。

\textbf{第三层——弦断裂与强子化（§4）：}
阐述静态夸克-反夸克势中的\textbf{弦断裂}（string breaking）现象——当夸克间距足够大，从真空中产生的动力学$q\bar{q}$对与原夸克对重组为两个静态介子，势在大距离处趋于平坦。在格点上，这需要包含动力学费米子（海夸克）的full QCD模拟——淬火近似下此现象被禁止。

\textbf{第四层——弱衰变与有效Hamiltonian（§5）：}
从电弱标准模型出发，W玻色子和Z玻色子交换在低能标下被积掉，产生\textbf{有效弱Hamiltonian}：
\begin{equation*}
H_{\text{eff}} = \frac{G_F}{\sqrt{2}} \sum_i V_{\text{CKM}}^i \, C_i(\mu) \, Q_i(\mu).
\end{equation*}
详述格点QCD在计算四夸克算符$Q_i$的强子矩阵元中的作用。讨论$K\to\pi\pi$衰变的两个核心挑战：紫外问题（算符的重整化和混合）和红外问题（Maiani--Testa定理——欧几里得关联函数在大体积极限下由零动量态主导，无法直接给出衰变振幅）。阐述Lellouch--L\"uscher关系如何克服这一困难——将有限体积中的$K\to\pi\pi$矩阵元与无穷大体积中的物理衰变振幅联系起来。

\textbf{第五层——CKM矩阵元与味物理（§6）：}
阐述格点QCD如何通过计算以下四类矩阵元来约束CKM幺正三角形的六个参数：
(i) \textbf{轻子衰变常数}：$f_{D_{(s)}}, f_{B_{(s)}}, f_K, f_\pi$——$|V_{cd}|, |V_{cs}|, |V_{ub}|$的直接测量；
(ii) \textbf{半轻子形状因子}：$B\to\pi\ell\nu, B\to D^{(*)}\ell\nu$——$|V_{ub}|, |V_{cb}|$的精确确定；
(iii) \textbf{$K$介子混合}：$B_K$参数——对幺正三角形的$\bar\rho\text{--}\bar\eta$平面提供强约束；
(iv) \textbf{$K\to\pi\pi$和$\epsilon'/\epsilon$}——CP破坏的直接与间接测量。

\textbf{第六层——粲介子衰变常数的现代精度（§7）：}
以本库中的粲介子衰变常数计算~\cite{CharmedClover}为具体实例，展示在13个$N_f=2+1$味系综、5个格距（$a\in[0.04,0.105]$~fm）上的连续极限计算。涵盖$D, D_s, D^*, D_s^*$的赝标量和矢量衰变常数、P波粲偶素的衰变常数、以及纯轻子衰变$D_{(s)}^*\to\ell\nu_\ell$的实验测试预言。

全篇以``从简单的两点函数提取衰变常数，到复杂的L\"uscher有限体积散射理论，再到标准模型味物理的完整CKM约束''为逻辑主线，系统展示了格点QCD如何成为弱衰变物理和味物理中不可替代的第一性原理理论工具。
\end{abstract}

\tableofcontents
\newpage

%==========================================================================
\section{引言：衰变——格点QCD中最具挑战性的物理领域}
%==========================================================================

\subsection{衰变在粒子物理中的核心地位}

强子的衰变是粒子物理中最基本的实验可观测量之一。从$\pi^+ \to \mu^+ \nu_\mu$（由$F_\pi$和CKM矩阵元$V_{ud}$控制）到$B^0 \to \pi^- \ell^+ \nu_\ell$（由$B\to\pi$形状因子和$|V_{ub}|$控制），每一种衰变道都编码了关于QCD束缚态动力学和电弱标准模型参数的关键信息。格点QCD——作为QCD的唯一第一性原理非微扰求解方法——在衰变物理中扮演着\textbf{不可替代}的角色：它将实验上测量的分支比和衰变宽度转化为标准模型基本参数（CKM矩阵元、低能常数）的无模型依赖约束。

\subsection{衰变格点计算的三大理论框架}

与强子质量的格点提取（从两点关联函数$C(t) \sim e^{-m_H t}$中简单读取）不同，衰变过程需要处理三个根本性的理论难题，它们分别催生了三种互补的理论框架~\cite{GattringerLang2010}：

\begin{table}[H]
\centering
\caption{格点QCD中衰变道的三类处理框架}
\label{tab:decay_frameworks}
\small
\begin{tabular}{p{2.5cm}p{4cm}p{2.5cm}p{4cm}}
\toprule
\textbf{框架} & \textbf{处理的衰变类型} & \textbf{格点算符} & \textbf{核心挑战} \\
\midrule
\textbf{衰变常数} & 轻子衰变（$\pi\to\ell\nu$, $K\to\ell\nu$, $D_{(s)}\to\ell\nu$） & 两点函数、轴矢流 & 重整化常数$Z_A$ \\
\textbf{L\"uscher方法} & 强衰变（$\rho\to\pi\pi$, 共振态） & 多点函数、变分法 & 有限体积中的能谱映射 \\
\textbf{有效弱Hamiltonian} & 弱衰变（$K\to\pi\pi$, $B$衰变） & 四夸克算符、三点函数 & 算符混合、Maiani--Testa定理 \\
\bottomrule
\end{tabular}
\end{table}

\subsection{本库中的核心文献}

\begin{table}[H]
\centering
\caption{本库中衰变道相关的核心文献分类}
\label{tab:decay_papers}
\small
\begin{tabular}{lp{9cm}}
\toprule
类别 & 关键文献 \\
\midrule
\textbf{教材} &
Gattringer \& Lang: Ch.11（衰变常数与低能参数、PCAC关系、弱衰变与有效Hamiltonian、L\"uscher散射方法、$K\to\pi\pi$、CKM矩阵）；Gupta: Ch.14, 16, 19（衰变常数、散射、味物理的实践指南）\\
\textbf{粲介子衰变常数} &
\texttt{Charmed meson masses and decay constants in the continuum from the tadpole improved clover ensembles.pdf}~\cite{CharmedClover}（13系综、5格距、$f_{D_{(s)}}$和$f_{D_{(s)}^*}$的连续极限——迄今最精确的粲衰变常数格点计算）\\
\textbf{重味/散射} &
\texttt{Kinematically-enhanced interpolating operators for boosted hadrons.pdf}（boost态内插算符对$\pi\pi$散射和重介子衰变的应用）；\texttt{Distillation at High-Momentum.pdf}（高动量蒸馏方法对衰变道中激发态控制的应用）\\
\textbf{夸克质量/低能常数} &
\texttt{Quark masses and low energy constants in the continuum from the tadpole improved clover ensembles.pdf}（与衰变常数密切相关的夸克质量与ChPT低能常数）\\
\bottomrule
\end{tabular}
\end{table}

%==========================================================================
\section{衰变常数：从两点函数到CKM矩阵元的约束}
%==========================================================================

\subsection{PCAC关系与$\pi$介子衰变常数}

\newcommand{\gtl}{Gattringer \& Lang}

$\pi$介子衰变常数$F_\pi$是强子物理中最重要的低能常数之一。它通过PCAC（部分守恒轴矢流）关系定义~\cite{GattringerLang2010}：

\begin{definition}[$\pi$介子衰变常数~\cite{GattringerLang2010}]
\begin{equation}
\boxed{\langle 0 | A_\mu^{a}(0) | \pi^b(P) \rangle = i \delta^{ab} F_\pi P_\mu},
\label{eq:Fpi_def}
\end{equation}
其中$A_\mu^a = \bar\psi \gamma_\mu\gamma_5 (\tau^a/2)\psi$是轴矢流，$|\pi^b(P)\rangle$是动量为$P$的$\pi$介子态。实验值$F_\pi = 92.4(3)$~MeV（本报告采用的归一化）。
\end{definition}

\textbf{格点提取：}从轴矢流-轴矢流关联函数的大$t$行为~\cite{GattringerLang2010}：
\begin{equation}
\langle A_4^+(p=0,t) A_4^-(0) \rangle \;\sim\; \frac{M_\pi F_\pi^2}{2}\, e^{-M_\pi t}, \qquad t \gg 0.
\label{eq:A4A4_correlator}
\end{equation}
通过关联$\chi^2$拟合提取指数前的系数$M_\pi F_\pi^2/2$，再结合从赝标量关联函数中独立确定的$M_\pi$，即可得到$F_\pi$。

\textbf{也常用赝标量-轴矢流交叉关联函数~\cite{GattringerLang2010}：}
\begin{equation}
\langle P^+(p=0,t) P^-(0) \rangle \sim \frac{M_\pi^3 F_\pi^2}{4 m_q^{(r)2} Z_P^2}\, e^{-M_\pi t},
\label{eq:PP_correlator}
\end{equation}
其中$m_q^{(r)}$是重整化的夸克质量。将式(\ref{eq:A4A4_correlator})和(\ref{eq:PP_correlator})结合，可从$F_\pi = 2Z_P m_q^{(r)} \sqrt{A/M_\pi^3} = 2Z_A m_{\text{AWI}}\sqrt{A/M_\pi^3}$中提取$F_\pi$。两种方法提供了冗余的交叉检验。

\subsection{$F_K$、$F_K/F_\pi$与$f_{D_{(s)}}$}

$K$介子衰变常数$F_K$的格点计算与$F_\pi$完全类比——仅需将轴矢流中的轻夸克替换为奇异夸克。$F_K/F_\pi$的比值是检验格点QCD味物理精度的基准可观测量——FLAG（Flavour Lattice Averaging Group）的世界平均值为$F_K/F_\pi = 1.193(3)$，与从$|V_{us}|$的实验约束一致。

对于粲介子（$D$和$D_s$），衰变常数通过类似的赝标量密度的真空-介子矩阵元定义~\cite{CharmedClover}：
\begin{equation}
\boxed{\langle 0 | \bar q \gamma_\mu\gamma_5 c | D_{(s)}(P) \rangle = i f_{D_{(s)}} P_\mu},
\label{eq:fD_def}
\end{equation}
其中$q = d$对$D^+$，$q = s$对$D_s^+$。重夸克（粲）的引入带来了额外的挑战：$am_c \sim 0.3-0.7$在典型格距上不可忽略，离散化误差需要通过高质量的重夸克作用量（如Fermilab作用量或相对论性重夸克作用量）来控制。

\subsection{衰变常数在CKM物理中的核心角色}

纯轻子衰变$\pi^+ \to \mu^+\nu_\mu$的宽度为：
\begin{equation}
\Gamma(\pi^+ \to \mu^+\nu_\mu) = \frac{G_F^2}{8\pi} |V_{ud}|^2 F_\pi^2 m_\mu^2 m_\pi \left(1 - \frac{m_\mu^2}{m_\pi^2}\right)^2,
\label{eq:pion_width}
\end{equation}
其中$G_F$是Fermi常数（从$\mu$子寿命精确已知）。因此，\textbf{格点QCD的$F_\pi$预言与轻子衰变分支比的实验测量相结合，直接确定了$|V_{ud}|$}——这是CKM矩阵元的最精确确定之一。

类似地，$D_{(s)}$介子的轻子衰变宽度为~\cite{CharmedClover}：
\begin{equation}
\Gamma(D_{(s)}^+ \to \ell^+\nu_\ell) = \frac{G_F^2}{8\pi} |V_{cd(cs)}|^2 f_{D_{(s)}}^2 m_\ell^2 m_{D_{(s)}} \left(1 - \frac{m_\ell^2}{m_{D_{(s)}}^2}\right)^2,
\label{eq:D_width}
\end{equation}
以及与对应矢量介子$D_{(s)}^*$的纯轻子衰变$D_{(s)}^* \to \ell\nu_\ell$（由衰变常数$f_{D_{(s)}^*}$控制）。BESIII合作组最近报告的$D^{*+}\to e^+\nu_e$的首个实验证据使格点QCD的$f_{D^*}$和$f_{D_s^*}$预言具有了直接的实验检验能力~\cite{CharmedClover}。

\subsection{Tadpole改进Clover系综上的粲衰变常数连续极限}

本库中的粲介子衰变常数计算~\cite{CharmedClover}代表了当前格点QCD在这一领域的最高精度。在13个$N_f=2+1$味系综、5个格距（$a\in[0.04,0.105]$~fm）上，CLQCD合作组完成了对所有S波开粲介子和P波粲偶素的衰变常数的连续极限外推。基于Tadpole改进的Clover费米子作用量，离散化误差在单个粗格距（$a\sim 0.1$~fm）处已压低至百分之几的水平。

\textbf{主要结果（部分）~\cite{CharmedClover}：}
\begin{itemize}[nosep]
    \item $f_{D_s} = 249.9(0.5)$~MeV（连续极限）——迄今最精确的$D_s$衰变常数格点结果；
    \item $f_D = 212.0(0.7)$~MeV；
    \item $f_{D_s^*}/f_{D_s} = 1.12(1)$——矢量与赝标量衰变常数的比值，对重夸克自旋对称性的检验；
    \item 结合PDG的轻子衰变分支比，$|V_{cs}| = 0.983(10)(13)(1)$和$|V_{cd}| = 0.220(2)(2)(6)$。
\end{itemize}

%==========================================================================
\section{强子衰变与L\"uscher有限体积方法}
%==========================================================================

\subsection{欧几里得时空中衰变振幅的不可直接获得性}

衰变是Minkowski时空中的实时动力学过程——末态粒子向无穷远处自由传播，其能谱是连续的。在欧几里得格点上（通过Wick旋转$t\to -i\tau$），连续能谱坍塌为离散的有限体积本征态——衰变振幅不能从格点关联函数中直接提取。这一根本困难由Maiani和Testa在1990年严格阐明~\cite{GattringerLang2010}：

\begin{quote}
``The Maiani--Testa theorem states that in Euclidean calculations in very large spatial volumes the matrix element is dominated by the pair of zero momentum single-particle states. In infinite volume in Euclidean correlation functions one obtains the average of matrix elements into in- and out two-pion states.''~\cite{GattringerLang2010}
\end{quote}

\subsection{L\"uscher公式：从有限体积能级到散射相移}

L\"uscher在1986--1991年的一系列奠基性论文中~\cite{Luscher1986}建立了将有限体积中的\textbf{离散双粒子能级}映射到无穷大体积中的\textbf{连续散射相移}的精确理论框架。核心物理图像（Gattringer \& Lang，§11.3.2~\cite{GattringerLang2010}）：

\begin{definition}[L\"uscher公式——一维简化推导]
考虑两个玻色子在长度为$L$的周期边界一维空间中散射。假设相互作用范围远小于$L$，则相互作用区域外的波函数为自由平面波，仅携带一个动量依赖的相移$\delta(k)$。周期边界条件要求：
\begin{equation}
e^{i k L + 2i\delta(k)} = 1 \quad\Longrightarrow\quad k_n L + 2\delta(k_n) = 2\pi n, \quad n\in\mathbb{Z}.
\label{eq:luscher_1D}
\end{equation}
对于给定的$L$，格点上可测量的是离散能级$W_n = 2\sqrt{m^2 + k_n^2}$。从不同$L$处的$W_n(L)$出发，通过式(\ref{eq:luscher_1D})可逐点重建相移$\delta(k)$。反之，如果$\delta(k)$的形式已知（如共振态参数化），有限体积能谱的大小依赖性可以被预言并与格点数据比较。
\end{definition}

\textbf{四维推广：}在实际QCD中，L\"uscher公式变得更为复杂~\cite{GattringerLang2010}：
\begin{equation}
\boxed{\delta(k) = \phi(q) \!\! \mod \pi, \qquad \tan(-\phi(q)) = \frac{q\pi^{3/2}}{Z_{00}(1; q^2)},}
\label{eq:luscher_4D}
\end{equation}
其中$q = kL/(2\pi)$，$Z_{00}$是推广的Riemann zeta函数。

\subsection{双粒子基态与$s$波散射长度$^{}$}

在阈值附近（$k\to 0$），散射由$s$波散射长度$a_0 = \lim_{k\to 0}\delta_0(k)/k$主导。L\"uscher公式在阈值展开下给出~\cite{GattringerLang2010}：

\begin{definition}[双粒子基态能量的体积依赖性~\cite{GattringerLang2010}]
\begin{equation}
\boxed{W = 2M_\pi - \frac{4\pi a_0}{M_\pi L^3} \left[ 1 + c_1\frac{a_0}{L} + c_2\frac{a_0^2}{L^2} + \mathcal{O}(L^{-3}) \right]},
\label{eq:threshold_expansion}
\end{equation}
其中$c_1 = -2.837297$和$c_2 = 6.375183$是仅依赖于格点几何形状的纯数值常数。
\end{definition}

\textbf{格点上的操作：}
\begin{enumerate}[nosep]
    \item 在2--4个不同体积$L$的系综上计算双$\pi$系统的基态能量$W(L)$（使用变分法从相关矩阵中提取）；
    \item 将$W(L)$对$1/L^3$标度拟合到式(\ref{eq:threshold_expansion})——提取散射长度$a_0$；
    \item 对非零总动量的推广可用于探测更高分波（$p$波、$d$波等）的相移。
\end{enumerate}

\subsection{$\rho\to\pi\pi$衰变：共振态格点计算的典范}

$\rho$介子是$\pi\pi$散射的$p$波共振态——它是强衰变过程的格点QCD计算的最成功案例~\cite{GattringerLang2010}。在$p$波中，L\"uscher公式通过非静止框架的推广（Rummukainen和Gottlieb, 1995）来处理非零总动量的双$\pi$系统~\cite{GattringerLang2010}：

\begin{enumerate}[nosep]
    \item 在格点上计算具有$\rho$介子量子数（$I=1$, $J^{PC}=1^{--}$）的双$\pi$关联函数的变分矩阵，提取若干个低激发能级$W_n(L)$；
    \item 利用推广的L\"uscher公式将每个$W_n$映射到一个动量$k_n$和相应的相移$\delta(k_n)$；
    \item 将提取的相移拟合到Breit--Wigner共振形式：
    \begin{equation}
    \frac{k^3}{W} \cot\delta_1(k) = \frac{6\pi}{g_{\rho\pi\pi}^2} (m_\rho^2 - W^2),
    \label{eq:bw_resonance}
    \end{equation}
    提取共振质量$m_\rho$（$\delta = \pi/2$处）和$\rho\pi\pi$耦合常数$g_{\rho\pi\pi}$——等价于衰变宽度$\Gamma_\rho$。
\end{enumerate}

\textbf{Gattringer \& Lang~\cite{GattringerLang2010}：}``In [34] a simple 2D system was studied which couples a heavier and two lighter bosons on the lattice, with mass and coupling parameters allowing for a decay like in the $\rho \to \pi\pi$ system. Figure 11.3 demonstrates the expected phenomenon of level-crossing avoidance, which leads to a resonating phase shift.''

在全QCD的动力学模拟中，CP-PACS合作组和HadSpec合作组已成功地在格点上完成了$\rho\to\pi\pi$衰变的计算，验证了L\"uscher架构对共振态物理的实用性。

\subsection{L\"uscher方法的局限与条件}

\begin{itemize}[nosep]
    \item 仅在\textbf{非弹性阈值以下}适用——当高于双$\pi$产生阈的通道（如$4\pi$, $K\bar K$）开启时，L\"uscher公式需要推广到耦合道形式；
    \item 需要$m_\pi L \gtrsim 4-5$以确保有限体积效应被指数压低；
    \item 对于宽共振（$\Gamma \gg \Delta E$，其中$\Delta E = 2\pi/L$是相邻能级的间隔），有限体积能级之间的回避交叉变得平滑，相移的提取精度恶化；
    \item 对于$\sigma/f_0(500)$等极宽的标量共振，L\"uscher方法面临严峻挑战，需要替代的分析策略。
\end{itemize}

%==========================================================================
\section{弦断裂与强子化}
%==========================================================================

\subsection{静态夸克-反夸克势中的弦断裂}

纯规范理论（淬火QCD）中，静态夸克-反夸克势在远距离处保持线性增长$V(r) \sim \sigma r$——这是禁闭的经典信号。然而，当\textbf{动力学费米子（海夸克）被纳入}时，介观图像发生质变~\cite{GattringerLang2010}：

\begin{center}
\fbox{\begin{minipage}{0.92\textwidth}
\small
\textbf{弦断裂（String Breaking）的物理过程~\cite{GattringerLang2010}：}

当夸克-反夸克间距$r$足够大（$V(r) \gtrsim 2M_B$，其中$M_B$是最轻重子-反重子对的质量），从QCD真空中产生一个动力学$q\bar{q}$对在能量上变得比继续拉长色弦更为有利。动力学$q\bar{q}$对中的$q$与原反夸克$\bar Q$重组为一个重-轻介子$(Q\bar q)$，而$\bar q$与原夸克$Q$重组为另一个重-轻介子$(\bar Q q)$。两个重-轻介子不再被色力束缚——它们自由分离，静态势在大$r$处趋于平坦（饱和到$\sim 2M_B$）。
\end{minipage}}
\end{center}

\textbf{格点验证：}在full QCD（$N_f=2$或$N_f=2+1$动力学夸克）的模拟中，弦断裂已通过静态夸克-反夸克势的``拉平''行为被数值验证。Gattringer与Lang指出~\cite{GattringerLang2010}：``This phenomenon is called string breaking and can be studied on the lattice.'' 需要强调的是，在淬火近似中弦断裂被\textbf{禁止}——海夸克环（闭合夸克线）被手动设定为1，意味着夸克-反夸克对不能从真空中自发产生。

\subsection{弦断裂与衰变常数的深层联系}

弦断裂不仅是禁闭动力学的一个有趣特征——它直接联系到\emph{衰变过程的格点可行性}。在full QCD中，动力学海夸克环使得强子可以衰变——例如$\rho \to \pi\pi$的末态$\pi$介子可以通过海夸克环的形成从真空中``物化''。在淬火近似中，这种衰变道被禁止——淬火近似是一个``稳定强子''的理论（忽略所有衰变）。这使得对衰变过程的格点研究\emph{本质上是全QCD（动力学模拟）的专有领域}。

%==========================================================================
\section{弱衰变与有效弱Hamiltonian}
%==========================================================================

\subsection{从电弱标准模型到有效弱Hamiltonian}

在高能标（$\sim M_W \approx 80$~GeV），弱相互作用由$W^\pm$和$Z^0$玻色子的交换来中介。在低能标（$\sim 1-2$~GeV，即强子物理的典型标度），重玻色子被``积分掉''（integrating out），产生\textbf{有效弱Hamiltonian}——四夸克定域算符的线性组合~\cite{GattringerLang2010}：

\begin{definition}[有效弱Hamiltonian~\cite{GattringerLang2010}]
\begin{equation}
\boxed{H_{\text{eff}} = \frac{G_F}{\sqrt{2}} \sum_i V_{\text{CKM}}^i \, C_i(\mu) \, Q_i(\mu)},
\label{eq:eff_weak_H}
\end{equation}
其中$G_F$是Fermi常数，$V_{\text{CKM}}^i$是相关的CKM因子，$C_i(\mu)$是\textbf{Wilson系数}（短程微扰QCD效应——在标度$\mu$处可微扰计算），$Q_i(\mu)$是\textbf{定域四夸克算符}（长程非微扰QCD效应——必须在格点上计算其强子矩阵元$\langle h'|Q_i|h\rangle$）。
\end{definition}

\subsection{四夸克算符的重整化与混合}

对于非手征Dirac算符（Wilson/Clover费米子），四夸克算符在重整化过程中与具有不同手征性和更低维度的算符发生混合~\cite{GattringerLang2010}。以$K\text{--}\bar K$混合为例：算符$O_{LL} = (\bar s \gamma_\mu(1-\gamma_5) d)^2$与$O_{PP} = (\bar s \gamma_5 d)^2$相混合。由于$\langle K|O_{LL}|K\rangle$在手征极限下正比于$m_q$（幂次压低），而$\langle K|O_{PP}|K\rangle$保持有限，微小的混合系数即可导致物理结果的大偏移——需要极高的非微扰精度来分离手征性不同的算符。使用Ginsparg--Wilson费米子（如Overlap或Domain-Wall）——它们在手征极限下自动保护手征性——显著简化了算符混合问题。

\subsection{$K\to\pi\pi$衰变与$\epsilon'/\epsilon$}

$K\to\pi\pi$衰变是CP破坏的两大实验信号——直接CP破坏（$\epsilon'$）和间接CP破坏（$\epsilon$）——的核心过程。格点QCD对$\epsilon'/\epsilon$的计算面临着前文所述的\textbf{双重挑战}：

\begin{enumerate}[leftmargin=*]
    \item \textbf{紫外（UV）问题：}有效弱Hamiltonian中四夸克算符$\Delta S=1$的重整化和与低维度算符（如$sd$、$s\gamma_5 d$、$s\gamma_\mu\gamma_\nu G^{\mu\nu} d$）的幂次发散混合。这些混合系数在$a\to 0$时以$1/a^n$的幂次发散——必须通过非微扰减除来消除~\cite{GattringerLang2010}。

    \item \textbf{红外（IR）问题（Maiani--Testa定理）：}欧几里得关联函数在大体积极限下由零动量双$\pi$态主导——但物理的$K\to\pi\pi$衰变中$\pi$介子携带非零动量。在有限体积中，Lellouch和L\"uscher~\cite{LellouchLuscher2001}导出了将有限体积$K\to\pi\pi$矩阵元与无穷大体积衰变振幅联系起来的精确关系——这是$K\to\pi\pi$物理格点计算的理论基石。
\end{enumerate}

\textbf{Gattringer与Lang~\cite{GattringerLang2010}：}``Lellouch and L\"uscher derived a relation between the $K\to\pi\pi$ matrix elements in finite volume and the physical kaon-decay amplitudes. ...A simple expression then relates the transition amplitude in finite volume to the decay rates in infinite volume.''

\subsection{$B_K$参数与味道物理}

$K^0\text{--}\bar K^0$混合由$\Delta S=2$的四夸克算符$O_{LL} = (\bar s \gamma_\mu(1-\gamma_5) d)^2$控制。其矩阵元通常参数化为\textbf{$B_K$参数}（``bag parameter''）：
\begin{equation}
\langle \bar K^0 | O_{LL} | K^0 \rangle = \frac{8}{3} M_K^2 F_K^2 B_K(\mu).
\label{eq:BK_def}
\end{equation}
$B_K$的格点QCD确定是约束CKM幺正三角形的$\bar\rho\text{--}\bar\eta$平面的\textbf{核心要素}之一。RBC/UKQCD合作组（使用Domain-Wall费米子）和其他合作组的$B_K$计算已将$B_K$的误差压至$<2\%$水平——是格点味物理的最大成功之一。

%==========================================================================
\section{CKM矩阵元与味物理}
%==========================================================================

\subsection{CKM幺正三角形的六个参数的格点约束}

Cabibbo--Kobayashi--Maskawa（CKM）矩阵$V_{\text{CKM}}$是标准模型中描述夸克味混合的$3\times3$幺正矩阵。其四个独立参数可参数化为Wolfenstein参数$(\lambda, A, \bar\rho, \bar\eta)$。格点QCD通过计算以下\textbf{四个类别的非微扰矩阵元}来约束这些参数~\cite{GattringerLang2010}：

\begin{table}[H]
\centering
\caption{格点QCD对CKM矩阵的约束输入}
\label{tab:ckm_inputs}
\small
\begin{tabular}{p{2cm}p{4cm}p{2.5cm}p{4.5cm}}
\toprule
\textbf{类别} & \textbf{格点可观测量} & \textbf{CKM约束} & \textbf{实验输入} \\
\midrule
轻子衰变常数 & $F_K/F_\pi$, $f_{D_{(s)}}$, $f_{B_{(s)}}$ & $|V_{us}|$, $|V_{cd}|$, $|V_{cs}|$, $|V_{ub}|$ & 轻子衰变分支比 \\
半轻子形状因子 & $B\to\pi$, $B\to D^{(*)}$, $K\to\pi$ & $|V_{ub}|$, $|V_{cb}|$, $|V_{us}|$ & 半轻子微分宽度 \\
$K\text{--}\bar K$混合 & $B_K$ & $\bar\rho$, $\bar\eta$ & $\Delta M_K$, $\epsilon_K$ \\
$B_{(s)}\text{--}\bar B_{(s)}$混合 & $f_{B_{(s)}}\sqrt{B_{B_{(s)}}}$ & $|V_{td}|$, $|V_{ts}|$ & $\Delta M_{B_{(s)}}$ \\
\bottomrule
\end{tabular}
\end{table}

\subsection{粲衰变常数对CKM的约束：典例分析}

以本库中的粲介子衰变常数计算~\cite{CharmedClover}为例~\cite{CharmedClover}：格点QCD给出的$f_{D_s} = 249.9(0.5)$~MeV与BESIII实验测量的$D_s^+\to\mu^+\nu_\mu$轻子衰变分支比相结合，得到$|V_{cs}| = 0.983(10)(13)(1)$。将此与CKM幺正性条件$|V_{ud}|^2 + |V_{us}|^2 + |V_{ub}|^2 = 1$对比，检验了CKM框架在粲夸克扇区的自洽性。格点精度（$\sim 0.2\%$对$f_{D_s}$）使得这一检验具有真正的辨别力——任何显著偏离幺正性的信号都将指向超越标准模型的新物理。

%==========================================================================
\section{总结与展望}
%==========================================================================

基于本库核心教材（Gattringer与Lang~\cite{GattringerLang2010}——衰变的完整格点理论）和约15篇研究论文，本文系统解析了格点QCD中衰变道的理论与计算框架。

\subsection*{核心结论}

\begin{enumerate}[leftmargin=*]
    \item \textbf{衰变常数是连接格点QCD与电弱标准模型的最直接桥梁：}$F_\pi$、$F_K$、$f_{D_{(s)}}$和$f_{B_{(s)}}$从真空-介子两点关联函数的指数衰减中提取——不涉及复杂的三点函数或末态相互作用。纯轻子衰变宽度与衰变常数的平方成正比——格点QCD提供的衰变常数值与实验分支比结合直接确定相关的CKM矩阵元。粲介子衰变常数的连续极限计算~\cite{CharmedClover}（$f_{D_s}=249.9(0.5)$~MeV）代表了当前格点QCD在这一领域的精度前沿。

    \item \textbf{L\"uscher有限体积方法使强衰变（共振态物理）的格点计算成为可能：}通过将有限体积中的离散双粒子能级映射到无穷大体积中的连续散射相移，L\"uscher公式（及其推广）绕过了欧几里得时空中衰变振幅不可直接提取的根本困难。$\rho\to\pi\pi$衰变是该方法的典范成功案例。L\"uscher方法的适用范围——非弹性阈值以下、$m_\pi L \gtrsim 4$、以及中等宽度共振——定义了当前格点共振态物理的前沿。

    \item \textbf{有效弱Hamiltonian方法将多标度问题分解为微扰（Wilson系数）和非微扰（四夸克矩阵元）两部分：}$K\to\pi\pi$衰变和$\epsilon'/\epsilon$的格点计算面临着双重的紫外和红外挑战。Lellouch--L\"uscher关系提供了处理红外问题（有限体积中$K\to\pi\pi$矩阵元到无穷大体积衰变振幅的映射）的系统框架。手征费米子（Ginsparg--Wilson）显著简化了紫外问题中算符混合的处理。

    \item \textbf{CKM幺正三角形的格点约束涵盖了轻子衰变、半轻子衰变和中性介子混合三类过程：}$B_K$参数（格点QCD精度$<2\%$）、$f_{B_{(s)}}\sqrt{B_{B_{(s)}}}$、以及$B\to\pi/D^{(*)}$的半轻子形状因子三者联合确定了$\bar\rho\text{--}\bar\eta$平面上的一个闭合约束区域——与$|V_{ub}|/|V_{cb}|$的比率、$\epsilon_K$和$\Delta M_{B_{(s)}}$的实验测量共同构成了当前标准模型味物理的最严格检验。
\end{enumerate}

\textbf{展望：}随着Exascale计算资源的到来，物理$\pi$质量下的full QCD模拟将使得$\pi\pi$、$K\pi$、$\pi K\bar K$等多强子道的关联L\"uscher分析成为可能——届时共振态谱学（$\rho$, $K^*$, $f_0(500)$等）将从格点QCD获得第一性原理的宽度和质量。重味衰变（$B\to D^{(*)}\ell\nu$, $B\to\pi\ell\nu$）在物理$b$夸克质量处的高精度计算将为$|V_{cb}|$和$|V_{ub}|$的持久分歧提供独立裁决。

%==========================================================================
\begin{thebibliography}{99}

\bibitem{GattringerLang2010}
C.~Gattringer and C.~B.~Lang,
\textit{Quantum Chromodynamics on the Lattice: An Introductory Presentation},
Lect.\ Notes Phys.\ \textbf{788}, Springer (2010).
\\\textbf{[来源:} \texttt{books/Quantum Chromodynamics on the Lattice.pdf}\textbf{]}
\\Ch.~11.1: 流算符定义、PCAC关系、$F_\pi$和$F_K$衰变常数的两点函数提取（Eqs.\,(11.10)-(11.12), (11.58)-(11.60)）。Ch.~11.3: L\"uscher有限体积方法——散射相移与体积依赖性能级、$\rho\to\pi\pi$衰变。Ch.~11.4.2: CKM矩阵、有效弱Hamiltonian、$K\to\pi\pi$衰变与$\epsilon'/\epsilon$、$B_K$参数（Eqs.\,(11.62)-(11.75), (11.107)-(11.108)）。Ch.~3.4: 弦断裂的物理图像。

\bibitem{Gupta1997intro}
R.~Gupta,
``Introduction to Lattice QCD,''
arXiv:hep-lat/9807028 (1997).
\\\textbf{[来源:} \texttt{books/INTRODUCTION TO LATTICE QCD.pdf}\textbf{]}
\\Ch.~14 (Wilson圈与静态势), Ch.~16 (衰变常数的实践提取), Ch.~19 (味物理与CKM矩阵)。

\bibitem{PeskinSchroeder1995}
M.~E.~Peskin and D.~V.~Schroeder,
\textit{An Introduction to Quantum Field Theory},
Westview Press (1995).
\\\textbf{[来源:} \texttt{books/An Introduction to Quantum Field Theory.pdf}\textbf{]}
\\Part III: QCD、部分子模型、弱相互作用与CKM矩阵的理论背景。

\bibitem{CharmedClover}
(CLQCD),
``Charmed meson masses and decay constants in the continuum from the tadpole improved clover ensembles,''
(年份/预印本)。
\\\textbf{[来源:} \texttt{docs/Charmed meson masses and decay constants in the continuum from the tadpole improved clover ensembles.pdf}\textbf{]}
——\textbf{13个系综、5个格距上的粲介子衰变常数连续极限计算。}$f_{D_s}=249.9(0.5)$~MeV, $f_D=212.0(0.7)$~MeV, $f_{D_s^*}$, $f_{D^*}$以及P波粲偶素衰变常数。CKM矩阵元$|V_{cs}|$和$|V_{cd}|$的确定。

\bibitem{Luscher1986}
M.~L\"uscher,
``Volume dependence of the energy spectrum in massive quantum field theories. I. Stable particle states,''
Commun.\ Math.\ Phys.\ \textbf{104}, 177 (1986);
``II. Scattering states,'' Commun.\ Math.\ Phys.\ \textbf{105}, 153 (1986).
——L\"uscher有限体积方法的奠基性论文。

\bibitem{LellouchLuscher2001}
L.~Lellouch and M.~L\"uscher,
``Weak transition matrix elements from finite-volume correlation functions,''
Commun.\ Math.\ Phys.\ \textbf{219}, 31 (2001), arXiv:hep-lat/0003023.
——将有限体积$K\to\pi\pi$矩阵元与无穷大体积衰变振幅联系起来的精确关系（Lellouch--L\"uscher关系——$\epsilon'/\epsilon$格点计算的基石）。

\bibitem{MaianiTesta1990}
L.~Maiani and M.~Testa,
``Final state interactions from Euclidean correlation functions,''
Phys.\ Lett.\ B \textbf{245}, 585 (1990).
——Maiani--Testa定理：欧几里得关联函数在大体积极限下的行为。

\bibitem{QuarkMassesClover}
``Quark masses and low energy constants in the continuum from the tadpole improved clover ensembles,''
\\\textbf{[来源:} \texttt{docs/Quark masses and low energy constants in the continuum from the tadpole improved clover ensembles.pdf}\textbf{]}
——与衰变常数密切相关的夸克质量和ChPT低能常数的连续极限计算。

\bibitem{GaugeSmearing}
``Gauge-Invariant Smearing and Matrix Correlators using Wilson Fermions at beta=6.2,''
\\\textbf{[来源:} \texttt{docs/Gauge-Invariant Smearing and Matrix Correlators using Wilson Fermions at beta=6.2.pdf}\textbf{]}
——包含赝标量衰变常数$f_p$的格点提取示例。

\bibitem{KinematicallyEnhanced}
``Kinematically-enhanced interpolating operators for boosted hadrons,''
\\\textbf{[来源:} \texttt{docs/Kinematically-enhanced interpolating operators for boosted hadrons.pdf}\textbf{]}
——对$\pi\pi$散射和重介子衰变道的boost内插算符优化。

\bibitem{DistillationHM}
``Distillation at High-Momentum,''
\\\textbf{[来源:} \texttt{docs/Distillation at High-Momentum.pdf}\textbf{]}
——高动量蒸馏方法对衰变道中激发态控制的应用。

\bibitem{HeavyMesonLCDAs}
``Calculation of heavy meson light-cone distribution amplitudes from lattice QCD,''
\\\textbf{[来源:} \texttt{docs/Calculation of heavy meson light-cone distribution amplitudes from lattice QCD.pdf}\textbf{]}
——重介子LCDAs对$B\to V$形状因子和遍举衰变的输入。

\bibitem{NovelQuarkField}
``A novel quark-field creation operator construction for hadronic physics in lattice QCD,''
\\\textbf{[来源:} \texttt{docs/A novel quark-field creation operator construction for hadronic physics in lattice QCD.pdf}\textbf{]}
——蒸馏方法对衰变性质的间接约束。

\bibitem{NoteGluonPDFs}
``Note of gluon PDFs,'' 内部笔记。
\\\textbf{[来源:} \texttt{docs/Note of gluon PDFs.pdf}\textbf{]}和\texttt{文档/note\_of\_gluon\_PDFs.pdf}

\end{thebibliography}

\vspace{0.5cm}
\noindent\rule{\textwidth}{0.5pt}
\small
\textbf{交叉参考建议：}本报告应结合同一\texttt{补充/}目录下的以下文件阅读：
\begin{itemize}[nosep]
    \item \texttt{格点QCD中的形状因子.pdf}——半轻子形状因子与过渡矩阵元（与衰变常数互补）
    \item \texttt{格点QCD中的关联函数.pdf}——两点/三点函数与衰变常数的指数拟合
    \item \texttt{格点QCD中的费米子方案.pdf}——不同费米子方案对算符混合的影响
    \item \texttt{格点QCD中的重夸克有效理论.pdf}——重味衰变（$B$, $D$介子）的理论框架
    \item \texttt{格点QCD中的连通图与非连通图.pdf}——非连通图在味单态矩阵元中的角色
    \item \texttt{格点QCD中的外推.pdf}——有限体积外推与手征外推
    \item \texttt{格点QCD蒸馏方法解析.pdf}——蒸馏对多强子道关联函数的增强
    \item \texttt{格点QCD中的重整化.pdf}——四夸克算符的重整化与混合
\end{itemize}

\end{document}
